\documentclass[11pt]{article}

\usepackage[margin=1in]{geometry}
\usepackage{times}
\usepackage[T1]{fontenc}
\usepackage[utf8]{inputenc}
\usepackage{microtype}
\usepackage{setspace}
\usepackage{booktabs}
\usepackage{natbib}
\usepackage{hyperref}
\usepackage{parskip}
\usepackage{url}

\setstretch{1.15}
\setlength{\parindent}{0pt}
\setlength{\parskip}{8pt plus 1pt}

\bibliographystyle{chicago}

\title{\vspace{-1cm}Website Access Rules and Visibility in AI-Powered Search}
\author{Xueyang Liu \\ ESADE Business School (Universitat Ramon Llull) \\ \texttt{xueyang.liu@esade.edu}}
\date{April 2026 --- Research Proposal}

\begin{document}

\maketitle

\vspace{-0.5cm}

% ════════════════════════════════════════════════════════════════════════════
\section{Introduction}

How investors acquire information about firms affects prices, the cost of
capital, and managerial discipline. Investors do not absorb every disclosure
on their own. They pay costs to search for information and to process it, and
these costs depend on the technology available to them. When EDGAR moved
corporate filings online, the cost of accessing regulatory disclosures fell,
and outsiders produced more information about firms \citep{gao2020informativeness,
goldstein2023edgar}. New technologies for finding information about firms
can change which firms investors notice and which they overlook.

This paper studies one such shift. AI-powered search systems, such as ChatGPT
Search and Perplexity, have become a common starting point for investors who
want information about a firm. A user types a question and the system returns
an answer drawing on web sources it has retrieved. Retail investors, and
analysts covering small firms, increasingly use these tools instead of
traditional search engines. A firm that the system frequently cites is more
visible to these investors than a firm it rarely mentions.

What such a system can say about a firm depends on what it can read. Some
firms restrict the AI systems that crawl their websites, either deliberately
or by inheriting a website template. When an investor asks an AI assistant
about these firms, the firms appear less often in the response and at lower
positions among cited sources. This is a form of information friction that
does not show up in disclosure data. A firm can meet every disclosure rule
and still be hard for investors to find through this channel.

I construct a dataset that flags firms whose web policies block AI search
systems. In a controlled experiment, these firms are cited less often, and
at lower positions, in answers produced by ChatGPT Search, Perplexity, and
Claude. For investors who use AI search, lower visibility raises the cost of
acquiring information about these firms. I then test whether this higher
cost shows up in standard measures of the firm's information environment.

The paper relates to several literatures. Work on the technology of information
acquisition \citep{da2011search, drake2012investor, gao2020informativeness}
shows that investors respond to the cost of finding information, and we add a
margin of cost variation that does not depend on a firm's disclosure choices.
Research on corporate websites as a disclosure channel
\citep{boulland2025website} treats website content as the relevant variable.
I show that website configuration matters as well, even when content is held
fixed. I also connect to recent work on the rise of AI crawlers and the data
commons \citep{longpre2024consent, cui2025robots} by tracing one
configuration choice through to investor outcomes.

% ════════════════════════════════════════════════════════════════════════════
\section{Institutional setting}

Three technical objects govern what an AI search system can read on a firm's
website.

\textbf{The \texttt{robots.txt} file.} A file at the root of the firm's
domain (\texttt{example.com/robots.txt}) that specifies which automated
agents (``crawlers'') can access which paths of the site. It is governed by
RFC~9309, which defines per-crawler groups of \texttt{User-agent},
\texttt{Allow}, and \texttt{Disallow} directives, with longest-prefix
matching and \texttt{Allow}-wins-on-tie. The file is unilaterally set by the
firm and publicly observable.

\textbf{The XML sitemap.} A list of URLs that the firm declares as
discoverable by search engines, defined by the Sitemaps protocol jointly
endorsed by Google, Yahoo, and Microsoft.\footnote{\url{https://www.sitemaps.org/}}
A firm that lists a URL in its sitemap is publicly signaling that page as
intended for crawler discovery.

\textbf{AI search crawlers.} Automated agents operated by AI companies to
retrieve web pages for inclusion in AI-generated answers. I focus on the
three crawlers whose downstream products have the largest consumer market
share and the most transparent citation behavior: OAI-SearchBot (ChatGPT
Search), PerplexityBot (Perplexity AI), and Claude-SearchBot (Claude). By
each vendor's own documentation, a site blocked for one of these crawlers
``will not appear in search answers'' produced by the corresponding product.

The treatment of interest is the overlap between these three objects. A
firm's sitemap declares certain URLs as discoverable, while its
\texttt{robots.txt} simultaneously prevents AI search crawlers from reaching
those same URLs. The firm's declared content is therefore public in principle,
but the AI systems routing a growing share of investor queries cannot read it.

% ════════════════════════════════════════════════════════════════════════════
\section{Data and measurement}

\subsection{Sample construction}

I start from the universe of U.S.~public firms in Compustat Fundamentals
Annual, and I extend the sample to include firms that exited between 2021
and 2025 to address survivorship concerns. For each firm, I clean the
\texttt{weburl} field, extract a canonical host (preserving subdomains, since
\texttt{robots.txt} is a host-level protocol), and identify 4{,}350 unique
hosts across approximately 9{,}500 firms with observable websites. Duplicate
hosts (e.g., 358 iShares ETFs all pointing to \texttt{www.ishares.com})
receive identical robots-level treatment and are handled via firm-level
sample flags.

\subsection{Crawler permissions}

For each host I fetch \texttt{robots.txt} and parse it under RFC~9309, and
for each of sixty-two documented AI crawlers (three consumer search crawlers,
the remaining thirteen AI search crawlers documented by Known Agents, and
additional sets of training and user-triggered crawlers)\footnote{I use the
taxonomy published by \url{https://knownagents.com/} and restrict the main
analysis to consumer-facing AI search crawlers with measurable downstream
products.} I compute two bot-level indicators: (i) whether the bot is
root-blocked (\texttt{Disallow: /}); and (ii) whether the bot has a
bot-specific \texttt{User-agent} rule, as distinct from inheriting a wildcard
(\texttt{User-agent: *}) rule. This distinction matters because roughly
ninety-nine percent of the access rules that apply to AI crawlers in my
sample come from generic wildcard rules that are likely to have been written
without AI in mind.

\subsection{The main treatment: sitemap--robots overlap}

For each host, I fetch the XML sitemaps declared in \texttt{robots.txt}
(resolving sitemap indexes recursively up to two levels, with gzip support
and a cap of 200{,}000 URLs per host). I then evaluate each sitemap URL
against each search crawler's applicable rules. The main treatment variable,
\texttt{search\_sitemap\_block}, is equal to one if at least one sitemap URL
is blocked for at least one of the three consumer search crawlers, zero if
the firm's sitemap URLs are observed and none are blocked, and missing if
overlap cannot be evaluated (fetch error, or no sitemap data).

This measure captures a directly observable conflict: the firm has listed
pages for search-engine discovery, yet its own access rules simultaneously
prevent AI search crawlers from reaching those pages. Under this definition,
one hundred thirty-seven firms (4.9\%) of the 2{,}796 non-fund firms with
evaluable sitemap data are classified as treated, versus forty-five percent
under a looser ``any search bot has any rule'' definition. The sitemap
overlap definition is tighter because it requires the restricted pages to
lie within the firm's own declared content, not merely within a generic
administrative path.

\subsection{Validation of the measurement}

I audit the content of the sitemap URLs used by treated firms. Using a
keyword classification adapted from \citet{boulland2025website}, less than
0.1\% of the sitemap URLs in the treated sample are technical or
administrative pages (e.g., \texttt{/login}, \texttt{/admin}). The remainder
are public content categories including product and service pages, news and
press releases, investor relations, and ``about the company'' pages. The
overlap therefore captures restrictions on economically meaningful content,
not incidental restrictions on backend paths.

% ════════════════════════════════════════════════════════════════════════════
\section{Preliminary evidence: experimental audit}

I validate the mechanism through a controlled experimental audit of AI
search citation behavior. Starting from the non-fund, Tier-2-evaluable
sample, I construct thirty matched pairs by exact matching within
SIC-2 $\times$ size-tercile cells: each pair contains one treated firm and
one matched control. For each of the sixty firms I then issue five
standardized natural-language queries covering the main information needs an
investor might have (what the firm does, financials, investor relations,
leadership, recent news). The queries are issued to three AI search
systems: ChatGPT Search, Perplexity AI, and Claude. I record, for each
firm$\times$query$\times$platform observation, whether the firm's official
website appears among the inline citations, its rank among unique cited
domains, and the full list of cited sources.

The audit is automated through the platforms' APIs wherever possible. The
Claude audit uses Anthropic's \texttt{web\_search} tool via a browser-side
artifact interface, which executes grounded queries with structured citation
output. The Perplexity audit uses Perplexity's Sonar API. The ChatGPT audit
is run manually given the absence of a publicly available ChatGPT Search API
at the time of writing. Pooling across platforms, treated firms are cited in
AI-generated answers about eight to eleven percentage points less often than
matched controls, and appear at lower positions among cited sources when
they are cited at all. The direction is consistent across platforms, and the
magnitude is largest for OpenAI (which has the highest consumer market
share and the highest webmaster awareness of its crawler) and smallest for
Claude (both figures are lower).

% ════════════════════════════════════════════════════════════════════════════
\section{Identification strategy}

The experimental audit validates the mechanism but does not identify the
effect of blocking on the firm's broader information environment. For that I
plan a staggered difference-in-differences design using firm-level variation
in the timing of policy changes.

I extend the cross-sectional measure into a monthly panel for 2021--2025
using snapshots of \texttt{robots.txt} and corresponding sitemaps retrieved
from the Wayback Machine's CDX API. This produces a firm-month panel of
approximately 600{,}000 observations. Preliminary panel statistics identify
about 2{,}600 firms that experienced at least one change in their AI
search-blocking status between 2021 and 2025. These ``switcher'' firms
provide identifying variation.

I will estimate the effect of sitemap--robots overlap on standard measures
of the firm's information environment (analyst forecast dispersion, analyst
coverage, bid-ask spread, Amihud illiquidity) using the heterogeneous-effect
estimator of \citet{callaway2021difference}. The staggered design addresses
selection on time-invariant firm characteristics; matching on observable
firm fundamentals addresses selection on observables at baseline.

As a placebo test, I exploit a variable I construct in the same step of the
pipeline: blocking of traditional search crawlers (Googlebot, Bingbot,
Baiduspider). Since a firm that blocks Googlebot essentially removes itself
from Google Search and would therefore face non-trivial costs in other
channels, meaningful variation in traditional-search blocking is rare. If
the estimated effects were driven by something other than the AI search
channel, I would expect similar coefficients on traditional-search blocking;
finding zero effects there would be consistent with the AI-search
interpretation.

% ════════════════════════════════════════════════════════════════════════════
\section{Expected contributions and next steps}

The project contributes a new firm-level measure of AI search visibility
constructed from the technical configuration of corporate websites, rather
than from disclosure content. The measure captures a form of information
friction that does not appear in standard disclosure data and that firms may
inherit unintentionally from default website templates. It complements
existing work on corporate websites as a disclosure channel
\citep{boulland2025website} by showing that configuration matters independent
of content, and it adds to the literature on the technology of information
acquisition \citep{gao2020informativeness, goldstein2023edgar} a margin of
cost variation that does not arise from firms' disclosure decisions.

The data construction pipeline is complete for the March 2026 cross-section.
The monthly panel is expected to finish in May 2026, after which I will
estimate the staggered difference-in-differences specifications and audit
robustness to alternative bot sets and outcome variables. A first full draft
of the paper is targeted for the end of May 2026.

% ════════════════════════════════════════════════════════════════════════════
\begin{thebibliography}{99}

\bibitem[Boulland et al.(2025)]{boulland2025website}
Boulland, R., Bourveau, T., \& Breuer, M. (2025). Corporate websites: A new
measurement frontier in accounting research. \textit{Journal of Accounting
Research}, 63(2), 545--589.

\bibitem[Callaway and Sant'Anna(2021)]{callaway2021difference}
Callaway, B., \& Sant'Anna, P.~H.~C. (2021). Difference-in-differences with
multiple time periods. \textit{Journal of Econometrics}, 225(2), 200--230.

\bibitem[Cui et al.(2025)]{cui2025robots}
Cui, T., Yang, Z., Liu, Y., Wang, X., \& Zhao, B.~Y. (2025). Robots.txt and
the governance of large language model crawlers. In \textit{Proceedings of
the 2025 ACM SIGSAC Conference on Computer and Communications Security}
(CCS '25).

\bibitem[Da et al.(2011)]{da2011search}
Da, Z., Engelberg, J., \& Gao, P. (2011). In search of attention.
\textit{The Journal of Finance}, 66(5), 1461--1499.

\bibitem[Drake et al.(2012)]{drake2012investor}
Drake, M.~S., Roulstone, D.~T., \& Thornock, J.~R. (2012). Investor
information demand: Evidence from Google searches around earnings
announcements. \textit{Journal of Accounting Research}, 50(4), 1001--1040.

\bibitem[Gao and Huang(2020)]{gao2020informativeness}
Gao, M., \& Huang, J. (2020). Informing the market: The effect of modern
information technologies on information production. \textit{The Review of
Financial Studies}, 33(4), 1865--1898.

\bibitem[Goldstein et al.(2023)]{goldstein2023edgar}
Goldstein, I., Yang, S., \& Zuo, L. (2023). The real effects of modern
information technologies: Evidence from the EDGAR implementation.
\textit{Journal of Accounting Research}, 61(5), 1699--1733.

\bibitem[Longpre et al.(2024)]{longpre2024consent}
Longpre, S., Mahari, R., Lee, A., et al.\ (2024). Consent in crisis: The
rapid decline of the AI data commons. \textit{Advances in Neural Information
Processing Systems}, 37.

\end{thebibliography}

\end{document}
